Large networks with hundreds of billions of parameters are extremely difficult to retrain, yet we often need to adapt a model to a new task or continue pretraining under a limited compute budget.

\cite{aghajanyan2020intrinsic} show that pretrained language models have a low ``intrinsic dimension''.
The Lottery Ticket Hypothesis~\cite{frankle2019lottery} similarly suggests that overparameterization can be reduced.
Low-rank approximation of large matrices and tensors has a long history in numerical linear algebra~(see, e.g., \citet{lubich2013projectorsplitting}, \citet{sulz2023numerical}).

Inspired by this, \cite{hu2021lora} hypothesized that weight updates during adaptation also have low intrinsic rank and proposed LoRA.
ReLoRA~\cite{lialin2023stack} extends LoRA by periodically merging low-rank adapters into the base weights and reinitializing them, so that the total update can become higher-rank over many cycles.
However, ReLoRA does not explicitly prevent successive adapters from occupying overlapping subspaces.
We propose SuperReLoRa, which keeps the multi-cycle merge-and-reinit structure but enforces column-space orthogonality of new LoRA factors at each reinitialization.
To mitigate suboptimal hard resets when a cycle has not yet converged, we soften the merge with transition momentum that partially retains the previous adapters.
Empirically, we compare LoRA, ReLoRA, and SuperReLoRa under identical hyperparameters and compute.
